% !TeX root = ../report_project_2stage_OTA.tex

\section{Project Overview}

This project develops a classical two-stage CMOS OTA using the SKY130A
open-source PDK.

The intended application is a low-power analog block representative of
PMIC and voltage-reference circuits, where DC gain, low current
consumption, stability and supply rejection are prioritized over
maximum speed.

The design follows a specification-driven methodology. First-order
analytical calculations are used to establish a coherent initial design
before transistor sizing. An ideal architectural model is then used to
study the fundamental two-stage behavior and compensation strategy
before implementing the complete transistor-level circuit.

Transistor sizing reuses the SKY130A device characterization obtained in
the preceding \texttt{project\_carac\_sky130A} project, including
$g_m/I_D$, $I_D/W$, intrinsic gain, threshold-voltage and saturation-voltage
data.

The design flow used throughout the project is:

\begin{enumerate}
    \item define the OTA specifications,
    \item select the topology and bias-current architecture,
    \item derive a first-order analytical design,
    \item study the ideal two-stage architecture and compensation,
    \item size the SKY130A transistors using the characterized device data,
    \item verify and iterate at transistor level,
    \item characterize the final OTA over the relevant operating conditions.
\end{enumerate}
